\begin{figure}[htbp]
\centering
\begin{tikzpicture}[
    font=\footnotesize,
    >=Stealth,
    orbitP/.style={draw=blue!75!black, very thick},
    orbitH/.style={draw=red!75!black, very thick},
    directionP/.style={draw=blue!75!black, thick, -{Stealth[length=2.3mm]}},
    directionH/.style={draw=red!75!black, thick, -{Stealth[length=2.3mm]}},
    guide/.style={draw=black!30, dashed},
    callout/.style={
        rectangle, rounded corners=2pt, fill=white,
        inner sep=1.5pt, font=\scriptsize
    },
    titleP/.style={
        rectangle, rounded corners=4pt, draw=blue!70!black,
        fill=blue!8, minimum width=48mm, minimum height=9mm,
        font=\bfseries, text=blue!70!black
    },
    titleH/.style={
        rectangle, rounded corners=4pt, draw=red!70!black,
        fill=red!8, minimum width=48mm, minimum height=9mm,
        font=\bfseries, text=red!70!black
    },
    note/.style={align=center, font=\scriptsize}
]

\path[use as bounding box] (-7.15,-3.65) rectangle (7.15,3.55);

% Títulos de los paneles
\node[titleP] at (-3.65,3.05)
    {ÓRBITA PARABÓLICA: $e=1$};
\node[titleH] at (3.65,3.05)
    {ÓRBITA HIPERBÓLICA: $e>1$};

% Fondos sutiles y separador
\fill[blue!2, rounded corners=8pt] (-7.0,-2.55) rectangle (-0.25,2.55);
\fill[red!2, rounded corners=8pt] (0.25,-2.55) rectangle (7.0,2.55);
\draw[black!18, thick] (0,-2.72) -- (0,2.72);

% -----------------------------------------------------------------
% Panel parabólico
% -----------------------------------------------------------------
\shade[ball color=yellow!85, draw=orange!80!black, thick]
    (-4.20,0) circle (0.42);
\node[font=\scriptsize\bfseries, text=orange!70!black]
    at (-4.20,-0.68) {Sol: foco};

% Eje de simetría y directriz conceptual
\draw[guide] (-6.80,0) -- (-1.00,0);
\draw[blue!40!black, densely dashed]
    (-5.52,-2.15) -- (-5.52,2.15);
\node[callout, text=blue!60!black, rotate=90]
    at (-5.74,1.40) {directriz};

% Parábola abierta hacia la izquierda con perihelio a la derecha
\draw[orbitP]
    (-6.65,-2.15)
    .. controls (-4.95,-1.52) and (-2.05,-1.02) .. (-1.55,0)
    .. controls (-2.05,1.02) and (-4.95,1.52) .. (-6.65,2.15);

\draw[directionP]
    (-6.28,-2.01)
    .. controls (-5.55,-1.74) and (-4.72,-1.52) .. (-4.05,-1.30);
\draw[directionP]
    (-3.90,1.33)
    .. controls (-4.70,1.55) and (-5.52,1.78) .. (-6.28,2.01);

% Perihelio y distancia q
\fill[blue!75!black] (-1.55,0) circle (2.4pt);
\node[callout, text=blue!70!black, above right=1mm and -1mm]
    at (-1.55,0) {perihelio};
\draw[<->, blue!60!black, thick]
    (-4.20,-0.90) -- (-1.55,-0.90);
\node[callout, text=blue!70!black] at (-2.88,-0.90)
    {$q$: distancia perihélica};

\node[
    note, draw=blue!55!black, fill=blue!5,
    rounded corners=3pt, inner sep=3pt
] at (-3.85,-2.05)
    {Energía orbital específica $\varepsilon=0$\\
     velocidad asintótica $v_\infty=0$};

\node[callout, text=blue!70!black] at (-6.10,-1.32) {entrada};
\node[callout, text=blue!70!black] at (-6.10,1.32) {salida};

% -----------------------------------------------------------------
% Panel hiperbólico
% -----------------------------------------------------------------
\shade[ball color=yellow!85, draw=orange!80!black, thick]
    (3.57,0) circle (0.42);
\node[font=\scriptsize\bfseries, text=orange!70!black]
    at (3.57,-0.68) {Sol: foco};

% Hipérbola de centro (1.30,0), semiejes a=1.75 y b=1.45.
% Sus asíntotas son y=+-(b/a)(x-1.30), y el Sol ocupa el foco derecho.
\draw[red!55!black, dashed, thick]
    (0.55,-0.62) -- (4.10,2.32);
\draw[red!55!black, dashed, thick]
    (0.55,0.62) -- (4.10,-2.32);
\node[callout, text=red!65!black, rotate=31]
    at (1.72,0.38) {asíntota};
\node[callout, text=red!65!black, rotate=-31]
    at (1.72,-0.38) {asíntota};

% Rama derecha de la hipérbola; el vértice corresponde al perihelio.
\draw[orbitH]
    plot[smooth] coordinates {
        (4.33,2.05) (3.91,1.60) (3.57,1.20) (3.30,0.80)
        (3.12,0.40) (3.05,0)
        (3.12,-0.40) (3.30,-0.80) (3.57,-1.20)
        (3.91,-1.60) (4.33,-2.05)
    };

\draw[directionH]
    (4.18,1.88)
    .. controls (3.82,1.52) and (3.48,1.14) .. (3.28,0.78);
\draw[directionH]
    (3.28,-0.78)
    .. controls (3.48,-1.14) and (3.82,-1.52) .. (4.18,-1.88);

% Perihelio y distancia q
\fill[red!75!black] (3.05,0) circle (2.4pt);
\node[callout, text=red!70!black, above left=1mm and -1mm]
    at (3.05,0) {perihelio};
\draw[<->, red!60!black, thick]
    (3.05,-0.90) -- (3.57,-0.90);
\node[callout, text=red!70!black] at (3.31,-0.90) {$q$};

\node[
    note, draw=red!55!black, fill=red!5,
    rounded corners=3pt, inner sep=3pt
] at (3.85,-2.05)
    {Energía orbital específica $\varepsilon>0$\\
     velocidad asintótica $v_\infty>0$};

\node[callout, text=red!70!black] at (4.72,1.62) {entrada};
\node[callout, text=red!70!black] at (4.72,-1.62) {salida};

% Comparación de excentricidad
\draw[black!55, thick, -{Stealth[length=2.2mm]}]
    (-6.45,-3.08) -- (6.45,-3.08)
    node[right, font=\scriptsize] {excentricidad $e$};
\fill[blue!8] (-6.35,-3.33) rectangle (0,-2.83);
\fill[red!8] (0,-3.33) rectangle (6.35,-2.83);
\draw[black!65, thick] (0,-3.38) -- (0,-2.76);
\node[font=\scriptsize, text=blue!70!black]
    at (-3.20,-3.08) {$e<1$: órbitas ligadas};
\node[
    circle, fill=blue!75!black, draw=white,
    minimum size=5mm, font=\tiny\bfseries, text=white
] at (0,-3.08) {$1$};
\node[font=\scriptsize, text=red!70!black]
    at (3.20,-3.08) {$e>1$: órbitas no ligadas};

\end{tikzpicture}
\caption{Comparación geométrica entre una órbita parabólica y una órbita hiperbólica. La parábola representa el caso límite $e=1$, con energía orbital específica nula, mientras que la hipérbola posee $e>1$, energía positiva y asíntotas que describen una velocidad residual distinta de cero lejos del Sol.}
\end{figure}